Tres cantidades fundamentales cuando se habla de circuitos electronicos son la corriente (A), el voltaje (V) y la resistencia (R). Un instrumento que permite la medicion de estas cantidades usandolo en diferentes configuraciones es el galvanometro de D'Arsonval.
% grafico del galvanometro
El Galvanometro usa el angulo de equilibrio entre el torque magnetico y el torque elastico para marcar una medida de la corriente que pasa por su emboninado. El torque magnetico viene dado por ${\bf r}\cross{\bf F}_m$, donde el brazo ${\bf r}$ es $a/2\hat{\bf r}$ y la fuerza magnetica $2n(ib\hat{\bf z})\times(B_0\hat{\bf r})$, con lo que el torque magnetico es $(a/2)nib_B_0\hat{\bf z}$ y como el torque elastico es $-k\theta\hat{\bf z}$, con lo que el torque total es $[(a/2)nibB_0-k\theta]\hat{\bf z}$, con lo que el punto donde el torque se anula y llega a un punto de equilibrio es $$\theta=\frac{nAB_0}{k}i$$, esto es una relacion lineal entre el corriente y el angulo permitiendo medor corrientes con facilidad midiendo un angulo.
% grafico de esquema de galvanometro
Cuando se hace referencia a un Galvanometro real, hay que considerar su resistencia interna, la cual se puede respresenta como $R_G$ en serie con el galvanometro. Este instrumento naturalmente tiene una corriente maxima que puede marcar $I_g^{max}$, la manera en la que se puede variar esta escala es conectando una resistencia en paralelo en la siguiente configuracion
% grafico de esquema de amperimetro
de esta manera, si queremos que la marca maxima de nuestro galvanometro marque una $I^{max}$, la resistencia que debemos colocar para nuestro amperimetro $R_A$ sera $$R_A=\frac{R_gI^{max}}{I^{max}-I_G^{max}}$$
Ademas de amperimetro, se puede usar el galvanometro en la siguiente configuracion para usarse como voltimetro
% grafico de esquema de voltimetro
donde se agrega una resistencia en paralelo para ajustar la escala de nuestro voltimetro $R_V$, donde si queremos que la marca maxima de nuestro Galvanometro marque $V^{max}$, $R_V$ sera $$R_V=\frac{V-I_G^{max}R_G}{I_G^{max}}$$